% Adición acumulativa. Insertar después de la suspensión esturmiana existente.
% Conservar íntegro el texto y la figura fig:holonomia-esturmiana.
\subsection{Nonadic connection, holonomy and monodromic representation}
\label{mono-ampl:conexion}

The regional connection already constructed is now studied through three objects:
transport of the complete state, representation of its windings and
suspension of its itineraries. The APP--TRIT--TPK operations and the
initial catalogue retain the preceding definitions. The necessary distinction
here concerns what each representation retains of the
same history(see \ref{hist:seccion}).

At prolongation index \(k\), a state chart preserves
\[
 \widehat x_k=(B_k,C_k,p_k,c_k,\Sigma_k,W_k,g_k,m_k).
\]
Here \(B_k\) collects the three regional blocks, \(C_k\) is the cylinder,
\(p_k\) the prefix, \(c_k\) the carry register, \(\Sigma_k\) the boundary
signature and \(W_k\) the incidence information. The phase is
\(g_k=1+[(k-1)]_9\); \(m_k\) counts complete windings. The reduction
\(\beta_k=[m_k]_{12}\) preserves the dodecaphasic position, while the
integer \(m_k\) distinguishes its successive windings.

\begin{definition}[Transport of a nonadic winding]
\label{mono-ampl:def-vuelta}
Let \(\mathcal R_k\subseteq\widehat X_k\times\widehat X_{k+1}\)
be the TPK admissible-extension relations. The realization of
holonomy at level \(k\) is the relational composition
\[
 \Gamma_{9,k}=\mathcal R_{k+8}\circ\cdots\circ\mathcal R_k.
\]
For \(s\geq1\), its prolongation to \(s\) windings is
\[
 \Gamma_{9s,k}=
 \Gamma_{9,k+9(s-1)}\circ\cdots\circ\Gamma_{9,k}.
\]
\end{definition}

The relational formulation preserves admissible continuations. On
an individually specified history, each arrow leads to the state actually
prolonged; the isolated visible block may still admit more than one
continuation. This difference is precisely the role of memory and
boundary conditions in determining the trajectory.

\begin{proposition}[Iteration of return and accumulated memory]
\label{mono-ampl:vueltas}
In every admissible composition of \(s\) windings,
\[
 g_{k+9s}=g_k,\qquad m_{k+9s}=m_k+s,\qquad
 \beta_{k+9s}=\beta_k+s\pmod{12}.
\]
\end{proposition}
\begin{proof}
The one-winding rule preserves phase and increments the counter by
one. Applying it to the already transported state repeats exactly these
two operations. Induction gives the first two equalities; reducing the
second modulo \(12\) gives the third.
\end{proof}

The counter thus provides an integer witness of the distinction between
windings. After \(12\) windings the finite phase pair is restored, but
the counter increases by \(12\). The nonsplit extension
\eqref{eq:extension108} expresses composition of the calendar through its
cocycle; the trajectory also preserves the integer lift of this datum.
Preservation here means transport of previous contributions,
not immobility of every coordinate.

\subsubsection{Bilateral representation of the return index}

The shift \(T\) of \eqref{eq:weyl-relojes} admits an
expression in phase and winding coordinates. Number phases from \(0\) to
\(8\), by shifting the origin relative to \(g_k\), and
write \(k=9m+r\), with \(0\leq r<9\), on the integer line
covering the phase cycle. On
\(\mathcal H_{\rm ind}=\ell^2(\mathbb Z)\), the basis
\(|k\rangle=|r,m\rangle\) expresses that same unitary shift:
\[
 T|r,m\rangle=
 \begin{cases}
 |r+1,m\rangle,&r<8,\\
 |0,m+1\rangle,&r=8.
 \end{cases}
\]
The inverse decreases \(r\) by one when \(r>0\), and takes
\(|0,m\rangle\) to \(|8,m-1\rangle\). Passage through the phase origin
thus incorporates the exact quotient unit.

\Needspace{8\baselineskip}
\begin{proposition}[Representation of cycle monodromy]
\label{mono-ampl:representacion}
Let \(S=T^9\). Then
\[
 S|r,m\rangle=|r,m+1\rangle,\qquad
 \mathbb Z\longrightarrow\mathcal U(\mathcal H_{\rm ind}),
 \quad s\longmapsto S^s
\]
is a faithful unitary representation of the cycle's winding group.
\end{proposition}
\begin{proof}
Nine shifts preserve the remainder upon division by \(9\) and
increment the quotient. Powers satisfy
\(S^{s+t}=S^sS^t\), and \(S^s|r,m\rangle=|r,m+s\rangle\) distinguishes
every integer \(s\). A permutation of an orthonormal basis is unitary.
\end{proof}

This representation corresponds to the cover of the cycle, whose winding
group is \(\mathbb Z\), and represents phase together with its counter. The
TPK state additionally contains the trajectory coordinates already described.
The bilateral extension admits negative indices as coordinates of the
cover; a trajectory beginning at an initial condition uses
its admissible half-line. This provides an operator inverse
of the index without turning every extension relation into a bijection of the
complete state.

With \(D_M|k\rangle=\lfloor k/9\rfloor|k\rangle\), the identity
\([D_M,S]=S\) expresses memory increment on the dense domain of
finitely supported vectors. If \(\mathcal I|k\rangle=|-k\rangle\) and
\(P_0\) projects onto indices divisible by \(9\), one also obtains
\begin{equation}
 \mathcal I T\mathcal I=T^{-1},\qquad
 \mathcal I D_M\mathcal I=-D_M-(I-P_0).
 \label{mono-ampl:inversion-contador}
\end{equation}
Indeed, for \(k=9m+r\),
\(\lfloor-k/9\rfloor=-m\) if \(r=0\), and equals \(-m-1\) if \(r>0\).
Inversion thus preserves the boundary term of Euclidean division.

\subsection{Block clock and capacity clock}
\label{mono-ampl:relojes}

The index of a refinement map and the number of available decimal
positions obey related but different laws. To
express them, reserve \(n\) for the ternary-block index and define
\[
 \rho=\log_{1000}729=\log_{10}9,\qquad
 \delta=1-\rho,\qquad
 \mathcal C_n=\lfloor(n+1)\rho\rfloor,\quad n\geq0.
\]
The symbol \(\mathcal C_n\) denotes capacity here, not the dodecaphasic
register \(K\). At the canonical phase origin, the preceding mechanical
word gives
\[
 Z_n=\sum_{j=0}^{n}z_j(0)=\lfloor(n+1)\delta\rfloor.
\]

\begin{proposition}[Exact balance between capacity and vacancies]
\label{mono-ampl:balance-capacidad}
For \(n\geq0\) and, in the second identity, \(n\geq1\),
\[
 \mathcal C_n=n-Z_n,\qquad
 \mathcal C_n-\mathcal C_{n-1}=1-z_n(0).
\]
\end{proposition}
\begin{proof}
Irrationality of \(\delta\) implies
\(\lceil(n+1)\delta\rceil=\lfloor(n+1)\delta\rfloor+1\).
Then
\[
 \lfloor(n+1)(1-\delta)\rfloor
 =n+1-\lceil(n+1)\delta\rceil=n-Z_n.
\]
Subtracting two consecutive identities completes the proof.
\end{proof}

An event \(z_n=1\) adds a block to the history while maintaining capacity
\(\mathcal C_n\); an event \(z_n=0\) increments both counters. In
any interval, capacity increments and vacancies sum
exactly to its length. Reducing modulo \(q\), for \(q=9\) or
\(q=108\), gives
\[
 [\mathcal C_n]_q=[n]_q-[Z_n]_q.
\]
Uniform block phase and capacity phase are related through
integrated vacancy memory. In particular, a capacity hold
and a complete winding of the connection have distinct mechanisms,
although both may preserve a phase observation.

Monotone inversion of the clock gives a first-arrival time:
\[
 T_{\rm cap}(a)=\min\{n:\mathcal C_n=a\}
 =\left\lceil\frac a\rho\right\rceil-1,
 \qquad a\geq1.
\]
With \(\sigma=\rho^{-1}-1\), the time required to gain \(b\) units
of capacity is
\[
 T_{\rm cap}(a+b)-T_{\rm cap}(a)
 =b+\lceil(a+b)\sigma\rceil-\lceil a\sigma\rceil.
\]
The difference of ceilings takes one of the two integers adjacent to
\(b\sigma\). Thus, \(9\) or \(10\) blocks are required to gain
\(9\) capacity units, and \(113\) or \(114\) to gain \(108\).
These are inverse-clock times measured from a first arrival; the
winding counter is always interpreted in the index to which it belongs.

\subsection{Symbolic suspension and aperiodic time crystal}
\label{mono-ampl:cristal}

Rotation coding brings together phase periodicity, local recurrence
and global aperiodicity. Let \(\mathcal X_\delta\subseteq\{0,1\}^{\mathbb Z}\)
be the closure, in the product topology, of the translates of the
bilateral word \(z_n(\theta)\). Denote the shift of its sequences by
\(\sigma_{\rm sh}\). Preserving a nonadic phase, the step is
\[
 F(r,z)=([r+1]_9,\sigma_{\rm sh}z),\qquad
 (r,z)\in C_9\times\mathcal X_\delta.
\]

\begin{definition}[Symbolic aperiodic time crystal]
\label{mono-ampl:def-cristal}
The unit-roof suspension of the preceding system is the space
\[
 \mathcal S_\delta=
 \bigl((C_9\times\mathcal X_\delta)\times\mathbb R\bigr)/\sim,
 \qquad (y,t+1)\sim(Fy,t),
\]
with flow \(\Phi_s[y,t]=[y,t+s]\). This model of temporal order, together with its
vacancy coding and winding-counter representation, is called a
\emph{symbolic aperiodic time crystal}.
\end{definition}

The adjective temporal designates the action of the flow and its return section;
aperiodicity characterizes its itineraries. The frequency of ones in each
sequence is \(\delta\): the uniform discrepancy bound strictly less than one,
obtained by telescoping in any window, persists under closure.
A periodic sequence would have rational frequency. Thus,
\(\mathcal X_\delta\) contains no periodic orbits. The suspension
has no closed orbits either: a flow return would require an integer
time and periodicity of the itinerary.

Finite words do recur. Length-\(m\) itineraries
arise from a partition of the circle by \(m+1\) points; each
interval is visited recurrently by the irrational rotation. Even
when times are restricted to a fixed phase, the step \(9\delta\) remains
irrational. This is why every word appears in all
\(9\) phases, and explains the previously established complexities
\[
 p(m)=m+1,\qquad p_9(m)=9(m+1),\qquad p_3(m)=3(m+1).
\]
Recurrence of local configurations coexists with the absence of a
global period. Its topological entropy is zero because these complexities
grow linearly.

Realization through the observables \(V_9,W_\delta\) and shift
\(T\) expresses the same order through the covariance relations
\eqref{eq:weyl-relojes}. Uniform-clock frequencies belong to
\(\frac19\mathbb Z+\delta\mathbb Z\), modulo \(1\). The capacity
clock additionally preserves the relation
\([\mathcal C_n]_9=[n]_9-[Z_n]_9\); it is not replaced by uniform phase
when interpreting returns. The term \emph{time crystal} is thus
defined through a mathematical dynamical system. Its physical realization
requires specifying energy dynamics, its temporal observables
and stability, data distinct from the present symbolic construction.

\subsection{Mirror symmetry of the regions and dual realization}
\label{mono-ampl:regiones}

Mirror interchange acts on the regional block
\(B=(u^\pi,u^e,u^\varphi)\in(\mathbb F_3^6)^3\) by
\[
 \mathfrak cB=(u^e,u^\pi,u^\varphi).
\]
Its fixed axis and aggregate are, respectively, \(u^\varphi\) and
\(s=u^\pi+u^e\). This fixedness concerns the involution in a fibre;
both coordinates may evolve during nonadic transport.
With \(q=u^\pi+u^e+u^\varphi\),
\(a=u^\pi+u^e-u^\varphi\) and \(d=u^\pi-u^e\), one recovers
\[
 u^\varphi=2(q-a),\qquad s=q-u^\varphi,\qquad
 u^\pi=2(s+d),\qquad u^e=2(s-d).
\]
The equalities hold componentwise in \(\mathbb F_3\), where
\(2^{-1}=2\). The mirror preserves \(q,a,s\) and reverses \(d\); this
antisymmetric coordinate individualizes the allocation that the aggregate preserves
without distinguishing. Time-index inversion \(\mathcal I\) and
channel involution \(\mathfrak c\) therefore act on
different data.

Phase return has an explicit witness:
\[
 \begin{aligned}
 B_1&=(012222\mid121221\mid112202),\\
 B_{10}&=(101222\mid112221\mid112002).
 \end{aligned}
\]
Both positions have phase \(1\); the block, axis and memory have
changed. Likewise, synchronization of the censuses at phases \(8\) and
\(9\) annihilates their relative differences, while the diagonal component
changes from \((59,59,59)\) to \((43,43,43)\). Equality of a relative observable
does not exhaust the state determining it.

The dual realization preserves that architecture. In Section
\ref{exc:doble-lectura}, the same prefix determines nested cylinders
on one side and a sequence of encoded words with their frames on the other.
The first map passes to the limit by contraction of diameters; the second
satisfies \(r_{n,m}Q_n=Q_mp_{n,m}\) and defines \(Q_\infty\).
Figure \ref{mono-ampl:fig-regional} anticipates these two maps.
The subsequent incidence route constructs codes, lattices and the
Moonshine realization in their own domains; the automorphism group
acts on that realization, not on the digits through an implicit
identification.

% Figura acumulativa; destino figures/monodromia_regional_ampliacion.tex.
\begin{figure}[!htbp]
\centering
\begin{tikzpicture}[x=1cm,y=1cm,>=Latex,font=\small,
  flow/.style={->,line width=.55pt},
  ann/.style={align=center,font=\footnotesize}]
\node[font=\bfseries] at (6.9,9.2)
  {Joint nonadic connection and regional coordinates};
\begin{scope}[shift={(3.4,6.8)}]
 \foreach \k in {1,...,9}{
  \pgfmathsetmacro{\ang}{90-40*(\k-1)}
  \coordinate (p\k) at (\ang:1.52);
  \fill (p\k) circle (1.4pt);
  \node[fill=white,inner sep=1.1pt,font=\footnotesize]
    at (\ang:1.83) {\(\k\)};
 }
 \foreach \k in {1,...,8}{
  \pgfmathtruncatemacro{\next}{\k+1}
  \draw[flow,shorten >=3pt,shorten <=3pt] (p\k)--(p\next);
 }
 \draw[flow,shorten >=3pt,shorten <=3pt] (p9)--(p1);
 \node[ann] at (0,.1) {\(g_k\in C_9\)\\\(\Gamma_{9,k}\)};
 \node[ann] at (0,-2.23) {One turn: \(g\mapsto g\), \(m\mapsto m+1\)};
\end{scope}
\node[ann,text width=5.25cm] at (10.1,7.75)
 {\(B_k=(u_k^\pi,u_k^e,u_k^\varphi)\)\\
 regional block in a fibre};
\node at (8.55,6.5) {\(u^\pi\)};
\node at (11.65,6.5) {\(u^e\)};
\draw[<->,line width=.65pt] (8.95,6.5)--node[above]{\(\mathfrak c\)}(11.25,6.5);
\draw[densely dashed,gray] (10.1,5.72)--(10.1,7.23);
\node[fill=white,inner sep=2pt] at (10.1,5.6){\(u^\varphi\)};
\node[ann] at (10.1,4.6)
 {Axis \(u^\varphi\) and aggregate \(u^\pi+u^e\)\\
 invariant under \(\mathfrak c\)};
\draw[flow] (5.45,6.8)--(7.25,6.8);
\draw[gray!60,line width=.3pt] (.55,4.02)--(13.25,4.02);
\node[ann,text width=12cm] (hist) at (6.9,3.47)
 {Compatible history: prefixes, cylinders, orientation, quotients,
 incidence frames, and memory};
\coordinate (fork) at (6.9,2.85);
\draw[flow] (hist.south)--(fork);
\draw[flow] (fork)--(3.6,2.2);
\draw[flow] (fork)--(10.3,2.2);
\node[ann,text width=5.6cm] at (3.6,1.65)
 {Archimedean realization\\
 \(I_{n+1}\subseteq I_n,\quad |I_n|=729^{-n}\)\\
 limit of the regional prefixes};
\node[ann,text width=5.6cm] at (10.3,1.65)
 {Incidence realization\\
 \(r_{n,m}Q_n=Q_mp_{n,m}\)\\
 encoded words and compatible frames};
\node[ann,text width=5.6cm] at (3.6,.42)
 {\(\pi,\varphi,e\); dodecaphasic determination of \(\alpha\)};
\node[ann,text width=5.6cm] at (10.3,.42)
 {Witt--Golay--Leech--\(V^\natural\)\\
 automorphisms of the exceptional realization};
\draw[flow] (3.6,1.03)--(3.6,.78);
\draw[flow] (10.3,1.03)--(10.3,.78);
\end{tikzpicture}
\caption{\fontsize{9}{11}\selectfont Nine phases of a connection and two realizations of the compatible
history. The cycle represents the phase; each turn increments the counter.
The regional diagram represents the involution \(\mathfrak c\), which
interchanges closure and propagation and fixes the self-scale axis in a fibre.
This invariance does not require the axis to remain constant during transport.
The two lower maps use the same prefix: one determines
its numerical limit, and the other preserves words and incidence frames. The
maps of this second realization are specified in Section
\ref{exc:doble-lectura}; the exceptional route is developed afterwards.}
\label{mono-ampl:fig-regional}
\end{figure}


\subsubsection{Rectilinear representation of phases and return}
\label{mono-recta:desarrollo}

The representation on the segment \([0,10]\) preserves the positional
origin of APP and orders its nine interior marks. In this diagram,
the integers \(1,\ldots,9\) are phase representatives. The horizontal
position indicates the ordering of the charts; the upper labels
specify regional operations. Figure
\ref{mono-recta:figura} thus brings together the initial line, fixation of the
self-scaling axis, the specular distribution of the two channels, and
return with memory(see \ref{mono-ampl:conexion}).

% Recuperación de fig_nonadic_route_mini.tex de la síntesis académica.
% Los rótulos designan operaciones regionales; la suma se efectúa en F_3^6.
\begin{figure}[H]
\centering
\begin{tikzpicture}[x=1.12cm,y=1cm,font=\small,>=Latex,
  guide/.style={line width=.45pt,black!55}]
 \draw[line width=.8pt] (0,0)--(10,0);
 \foreach \k in {0,...,10}{
  \draw[guide] (\k,-.09)--(\k,.09);
  \node[below=2.5mm] at (\k,0) {$\k$};
 }
 \foreach \k in {1,...,9}{\fill (\k,0) circle (1.7pt);}
 \node[anchor=west,font=\footnotesize] at (0,3.05)
   {Generating segment $[0,10]$ and nine phase representatives};
 \draw[guide] (5,.12)--(5,1.7);
 \node[above,align=center,font=\footnotesize] at (5,1.7)
   {axis fixation\\$u^\varphi$};
 \draw[decorate,decoration={brace,amplitude=4pt},guide]
   (6,.65)--(8,.65);
 \node[above=5pt,align=center,font=\footnotesize] at (7,.65)
   {regional aggregate\\$u^e+u^\pi$};
 \draw[guide] (9,.12)--(9,1.7);
 \node[above,align=center,font=\footnotesize] at (9,1.7)
   {specular involution\\$\pi\leftrightarrow e$};
 \draw[->,line width=.7pt] (9,-.7)
   .. controls (9,-1.55) and (1,-1.55) .. (1,-.7);
 \node[fill=white,inner sep=3pt,font=\footnotesize] at (5,-1.23)
   {return $9\to1$;\quad $m\mapsto m+1$};
 \node[font=\footnotesize] at (5,-1.95)
   {$w_6\to w_{12}\to w_{18}\to w_{24}\to w_{30}\to R_{36}\to G_9$};
\end{tikzpicture}
\caption{Rectilinear representation of the joint nonadic connection. The mark
\(5\) locates the first strong saturation and the regional self-scaling
axis; the brace \(6\)--\(8\) identifies the ternary aggregate
fixed in each fibre; phase \(9\) selects the closure
orientation. Return preserves phase and advances memory. The
abscissae number operational positions, whereas the labels
\(u^\pi,u^e,u^\varphi\) designate blocks of \(\mathbb F_3^6\).}
\label{mono-recta:figura}
\end{figure}


Axis fixation is formulated in the fibre of a given boundary. If
\(q=u^\pi+u^e+u^\varphi\) and
\(a=u^\pi+u^e-u^\varphi\), then
\[
 u^\varphi=2(q-a),\qquad s=u^\pi+u^e=q-u^\varphi
 \quad\text{in }\mathbb F_3^6.
\]
The data \((q,a)\) determine the axis and the aggregate; the ordered distribution
of \(s\) retains the specular freedom resolved by prolongation.
The fifth phase has one local solution and one extension, with boundary
datum \(c=111111\) and residual signature \(\rho_W=000\). This is the
meaning of its first strong saturation. The symbol \(\varphi\)
above the mark \(5\) identifies the regional axis involved
in that operation.

The intermediate phases preserve this decomposition within each fibre,
while their boundary data evolve. The complete census of
local multiplicities and extensions of the distinguished branch is
\[
\begin{array}{c|rrrrrrrrr}
\text{phase}&1&2&3&4&5&6&7&8&9\\\hline
N_{\rm loc}&3&2&5&4&1&1&3&3&2\\
N_{\rm ext}&2&5&4&1&1&3&3&2&1
\end{array}
\]
Phase \(6\) combines local uniqueness with three prolongations; phase
\(7\) resolves a triple specular fibre; phase \(8\) synchronizes
the three regional censuses. Phase \(9\) retains two local
distributions with the same axis and aggregate, one of which
prolongs. These counts express different functions within the
same connection and justify the marks on the line.

In particular, for phases \(6,7,8\), the regional sums are
\(012100\), \(101001\), and \(112000\), respectively. The conservation
represented by the brace is invariance under specular distribution
within each fibre. Subsequent real evaluation applies the
readouts of compatible cylinders to complete histories; the
internal sum represented in the figure belongs to the ternary domain.

Finally, the arc \(9\to1\) represents passage from one turn to the
next. The formulas in Section
\ref{mono-ampl:conexion} preserve the phase remainder and increment the
turn counter. The rectilinear diagram and the cyclic diagram in
Figure \ref{mono-ampl:fig-regional} therefore express two complementary
aspects of the same holonomy: the ordering of positions and
composition of return on the state with memory.


\subsection{Temporal transport and tritic exponential composition}
\label{mono-ampl:euler}

The Euler relations in Section \ref{euler-trit-articulacion}
describe local realizations of the three regimes. For a fixed
signature \(\tau\), the generator satisfies \(J_\tau^2=-\tau I\); algebraic
conjugation takes \(J_\tau\) to \(-J_\tau\). Consequently,
\[
 \exp(sJ_\tau)\exp(tJ_\tau)=\exp((s+t)J_\tau),\qquad
 \overline{\exp(tJ_\tau)}=\exp(-tJ_\tau),
\]
and the product of an evolution with its conjugate is the unit. In the
elliptic type this is realized by \(\cos^2t+\sin^2t=1\); in the hyperbolic type,
by \(\cosh^2t-\sinh^2t=1\); in the parabolic type,
\((I+tN)(I-tN)=I\), since \(N^2=0\). The same conservation law
admits three quadratic expressions.

The elliptic realization closes a winding; the nilpotent one preserves a
linear displacement; the hyperbolic one separates expansion and contraction with
unit determinant. At the regional parameters already generated,
\[
 \exp(\pi J_{\mathrm e})=-I,\qquad
 \exp(J_{\mathrm p})=I+J_{\mathrm p},\qquad
 \exp(2\log\varphi\,J_{\mathrm h})=F_{\rm av}^{\,2}.
\]
The evaluation \(\exp(I)=eI\) recognizes the unit of the common exponential
law. The number \(e\) is not assigned exclusively to the parabolic regime.

Transport of a local realization by an isomorphism \(U\) preserves
these relations: if \(J'=UJU^{-1}\), the exponential series gives
\(U\exp(tJ)=\exp(tJ')U\). This identity transports the regime and its
norm. A transition between different quadratic types belongs,
by contrast, to the TPK selection of sheets and regimes; it is not a conjugation
changing the value of \(J^2\). The temporal product therefore preserves the
order of transitions and information about their endpoints.

This articulation distinguishes closure of a local representation from
holonomy of a history. One may have \(\exp(2\pi J_{\mathrm e})=I\)
at the same time as \(S|r,m\rangle=|r,m+1\rangle\). The first equality
restores regime orientation; the second records another winding. The
discrete structure preserves both: local return and the accumulated provenance
of the state realizing it.
